\section{Elementary implementation and a finite parameter choice}\label{s:compile}\label{sec:compilation}
The controller has a finite alphabet but an enormous local rule. We
first show that its size affects only on-site depth, and then choose
padding to make that depth compatible with the source estimates.

\subsection{Binary alphabets and a streamed Boolean circuit}
Encode a chart word in \(r=\lceil\log_2|\mathcal A|\rceil\) bits.
Unused codes mean the bottom symbol. Dephase every retained controller
word before using it, decode it, and write canonical output codes. This
makes bottom aliases harmless even for an arbitrary initial state.
For helper fields use seven bits: one low code, one unknown code and
126 high codes. Uniform whole-site replacement supplies uniform bits
in these registers. It leaves the needed upper and lower helper laws
\begin{equation}\label{eq:helper-law}
 u_d=\frac{127}{128}\bigl[1-(1-\mu)^{\ell_d}\bigr],\qquad
 l_W=\frac{126}{128}\bigl[1-(1-\mu)^W\bigr].
\end{equation}
Here \(\mu\) is the aggregate semantic refresh rate. The short experiment uses
only \(u_d\le\mu\ell_d\); the long experiment uses
\(126/128\ge .98\), so \(l_W>.9C\) in the parameter range below.

For a classical function \(f:\{0,1\}^m\to\{0,1\}^r\), start with the new output word and scratch zeroed, then enumerate
inputs \(w\). Negate wires corresponding to negative literals of \(w\),
compute their conjunction in zeroed scratch wires by Toffoli gates,
copy the resulting indicator into each output bit with \(f(w)_j=1\),
and uncompute. The cases \(m=0,1\) use constants and CNOT. A safe bound is
\[
 G_f\le2^m(8m+2r+20)+10(m+r+1).
\]
There are \(2^m\) rows and \(r2^m\) truth-table bits; neither need be
stored simultaneously. Computing \(f(w)\) from the finite semantic rule
and writing the corresponding gates is a terminating finite algorithm.
It uses no answer to a halting question. All its gates are on site.
The same construction computes finite matching rules for the quantum
processor. Scratch and transport registers reset before first use and
before the end of each round.

\subsection{Classical transport and the source name}
For each displacement and bit allocate a distinct sender and a distinct
zeroed receiver. Fan out the old dephased word locally using CNOTs.
In one tensor-disjoint layer send each sender exactly once into its
receiver by CNOT. No received word is subsequently forwarded to another
site. The receiver then participates only in local computation.

A flag before transport at a sender corrupts copies of that site's
old word. A faulty zero receiver affects only that destination's new
computation. A flag after transport likewise affects only the new
computation at its site. Consequently all microscopic flags at one
site in one round have one spatial source name, with old and new
appearances at the two adjacent formal times. Reusing the same received
word for all roles reading that old site is essential. A spatial SWAP
implemented by three noisy CNOTs would allow a middle fault to be
forwarded; the one-way transport supplies the attribution used here.

All information-moving quantum gates occur before the long on-site
Boolean calculation. The fixed native candidate-port
hardware has a bounded propagation cone. A flag during that hardware
is an arbitrary channel on a bounded semantic support. During the later
Boolean calculation quantum data is idle until the final erasure.
The spatial quantum-fault cone therefore does not grow with Boolean
depth or helper-lane length. Remote controller errors are charged by
the controller's derivation, rather than by a new quantum fault at
every induced erasure.

\subsection{Unconditional native operations on private ports}\label{d:dispatch}
For every native gate type in \eqref{d:native} and every translated
operand tuple, allocate private type and position ports. Reset them,
route matching semantic operands into their ports by metadata-controlled
local swaps, apply the native gate unconditionally to every translated
role tuple, reverse the routes, and reset the ports. On the clean
circuit, either all operands of a matching active tuple are routed or
none are. Swap conjugation implements the native channel on data and
reference; an inactive gate acts only on disposable ports.
A fault inside an inactive route can affect its data cell and is counted
as a local dirty location.

Position role \(j\) at site \(u\) belongs only to tuple base
\(u-\delta_j\). Therefore all translated gates are tensor-disjoint,
even on small tori where sites coincide. Serialize types, incompatible
routes and elementary decompositions. Malformed masks still define a
valid circuit on the same finite neighborhood. One semantic qubit per
chain cell suffices; code-sized work is distributed over additional
cells rather than added to this fixed role catalogue.

The controlled route is the Fredkin word
\(\mathrm{CNOT}(b;a),\mathrm{Toffoli}(c,a;b),\mathrm{CNOT}(b;a)\).
For \(c=0\) the CNOTs cancel; for \(c=1\) it swaps the two factors.
Also \(\mathrm{SWAP}=\mathrm{CNOT}_{ab}\mathrm{CNOT}_{ba}
\mathrm{CNOT}_{ab}\), \(\mathrm{CZ}=H_b\mathrm{CNOT}_{ab}H_b\),
\(Z=HXH\), and discard is \(R_0\) on its retained slot.
Consequently every physical instruction belongs exactly to the seven
gate inventory in Section~\ref{s:model}. Every sublayer in these words
receives noise. The real closed native subsystem of
Section~\ref{sec:quantum} supplies every intended opcode.

Let \(\mathcal M\) denote all autonomous classical fields. Initially dephase
\(\mathcal M\); its own gates are Boolean basis permutations, resets and
dephasings. In an interaction with quantum data, \(\mathcal M\) is only a
control and never a target, at every elementary sublayer. Depolarization
preserves block diagonality in \(\mathcal M\). Induction shows that its noisy
history law is independent of the quantum input even for malformed
initial words and intermediate faults. Quantum-derived records stay
in the quantum track. Old calendar symbols control the current port;
old independent copies compute the next classical symbols.

\subsection{A terminal erasure that cuts dependence}
For each data qubit \(a\), prepare a helper \(b\) by \(R_0,H,D_Z\),
forgetting the measurement outcome. Conditional on every intervening
flag history the result is exactly \(I/2\). With no flag this follows
from the three gates; after a flag the remaining operations are unital
on an already maximally mixed local quantum register.
Let \(c\) be the classical erase bit. The sequence
\[
 \mathrm{CNOT}(b\to a),\quad
 \mathrm{Toffoli}(c,a\to b),\quad
 \mathrm{CNOT}(b\to a)
\]
is identity on \(a,b\) when \(c=0\), and swaps them when \(c=1\).
Apply it to every retained quantum factor, including native syndrome,
ancilla and work factors, and then reset the helpers.
A clean block with \(c=1\) replaces the data by \(I/2\) and discards
its entire previous reference correlation. In \(\mathrm W\), a site flag within the
block replaces all local quantum factors; the remaining controlled
unitaries preserve their maximally mixed state. The clean replacement identity is a complete channel identity,
including on inputs entangled with a reference.

For arbitrary-input forgetting we use the direct dependence statement:
in every round, every path from retained input quantum factors to
retained output quantum factors crosses that round's final gatekeeper
block. A clean bottom trigger cuts the path, including paths through
temporary hardware. Only preparation, storage, noise and final
work discard/reset follow the block. This statement is stronger than
resetting helpers at endpoints: a swap into a helper followed by a
swap back after erasure would violate it. The actual schedule satisfies
the per-round condition.

\subsection{Whole-site realization \texorpdfstring{\(\mathrm W\)}{W}}
Let \(J\) be the number of internal layers, including their following
noise. Add \(P_{\rm pad}\) identity/noise layers at the end of each round. At
site \(z\) and round \(t\), let \(A_{z,t}\) mean at least one internal
replacement and \(E_{z,t}\) at least one padding replacement. Then
\begin{equation}\label{eq:macro-rates}
 \nu=1-(1-\eps)^J,\qquad \mu=1-(1-\eps)^{P_{\rm pad}},\qquad
 H=A\lor E,
 \qquad p=1-(1-\nu)(1-\mu)\le2\mu
\end{equation}
if \(P_{\rm pad}\ge J\). These are independent families across physical
site/round names; \(A\) and \(E\) are independent. Padding can be
represented by one replacement mark, the last one, when \(E=1\).
Constituent marks are independent uniform variables conditional on that
one site flag. Raw source tests may use \(H\); baseline helper tests use
only padding flags and marks.

For \(\mu\le1/2\),
\[
 \nu\le J\eps\le2J\mu/P_{\rm pad}.
\]
A source bin using at most \(h_0W\) site/round source names, with \(W\le C/\mu\),
is spoiled by internal flags with probability at most
\begin{equation}\label{eq:overlay}
 \alpha\le 2h_0JC/P_{\rm pad}.
\end{equation}
The overlap graph has fixed finite colour divisors \(D_{\rm bin}\)
and \(D_{\rm hist}\), chosen at least one above its maximum degree.
They depend on the spatial supports within one round, not on \(J,W,n\).
Independence of \(A\) and padding experiments, colouring, and
Cauchy--Schwarz give the augmented estimates
\begin{equation}\label{eq:gamma}
 \gamma_0\le\sqrt{2C}+\eta_B^{1/6}
       +\alpha^{1/(2D_{\rm bin})},\qquad
 \theta=\eta_B+\alpha^{1/D_{\rm hist}}.
\end{equation}
Distinct labels within one selector obey the joint bound \(\gamma_0^g\), as used in
the quantum section; \(\theta\) is the historical error weight. Labels,
rather than controller-induced erasures, carry this bound.

\subsection{Independent-qubit realization \texorpdfstring{\(\mathrm Q\)}{Q}}\label{s:probe}
The \(\mathrm Q\) realization supplies the same round interface by a separate
physical circuit. It adds a local sampling experiment and controlled
refresh of every semantic register, including all quantum, calendar,
selector and helper fields (and the profile in the folded appendix). Denote by \(N_{\rm np}\) the number of
non-probe wires per site: the core, one mixed helper per semantic wire,
two OR scratch bits and all fixed work. This count is independent of
the number \(M\) of probes.

To see this, compute OR serially. Given a basis-valued accumulator
\(a\), probe \(p\), and target \(b=0\), apply CNOT \(a\to b\),
CNOT \(p\to b\), and Toffoli \((a,p;b)\). Then
\(b=a\mathbin{\mathrm{XOR}}p\mathbin{\mathrm{XOR}}ap=a\lor p\).
Reset \(a\), copy \(b\to a\), and reset \(b\). Repeating this
six-gate pattern uses two scratch wires.

After the core and its terminal quantum cut, reset all \(M\) probes to
zero. Including the noise following that reset as the first layer,
let them undergo \(K_{\rm samp}\) sampling noise layers, with identities in the
remaining layers. Every non-probe wire idles and is noisy as well.
At the endpoint, each probe has independent probability
\(q_{\rm end}=[1-(1-\eps)^{K_{\rm samp}}]/2\) of value one. All its sampling states are
diagonal. Define the exogenous event \(E\) as the OR of these endpoint
values, before any subsequent noisy OR computation. Its exact rate is
\begin{equation}\label{eq:q-refresh}
 \mu(\eps)=1-\left[\frac{1+(1-\eps)^{K_{\rm samp}}}{2}\right]^M.
\end{equation}
It is independent of the semantic input and reference because of the
initial probe reset.

Compute OR, prepare private mixed helpers by \(R_0,H,D_Z\), and use
the controlled-swap gadget to refresh every semantic factor when the
computed control is one. Reset all swapped-out contents, probes and
scratch before declaring the round's outputs. No discarded quantum
state is recovered afterward. Let \(J\) count all non-sampling layers,
including the core and this finite OR/refresh/discard circuit. It is
fixed after \(M\), before \(K_{\rm samp}\).

Define \(A\) as the union of all flags on non-probe wires during
sampling and all flags on every wire in the non-sampling layers.
The primitive coins defining \(A\) are disjoint from all probe
sampling flags and endpoint marks defining \(E\). Hence \(A,E\) are
exactly independent iid site/round fields, with
\begin{equation}\label{eq:q-overlay}
 \nu\le[N_{\rm np}K_{\rm samp}+(M+N_{\rm np})J]\eps.
\end{equation}
Post-sampling probe faults, including an incorrect OR, belong to
\(A\). On \(E\land\neg A\) the semantic output is exactly the
independent product of \(I/2\) on every factor; on
\(\neg E\land\neg A\) it equals the clean core output. We use only
this clean refresh implication in \(\mathrm Q\). Actual output on \(A\) may be
arbitrary; in particular, \(E\) alone does not imply uniform output.
Counterfactual baseline mixed-helper marks can be taken independent and
uniform on dirty rounds, agreeing with the actual refresh on clean
rounds. These marks support the original helper experiments, while
\(A\) enlarges every dirty experiment to an intrinsic exception.

For \(MK_{\rm samp}\eps\le1/2\), elementary union and product bounds give
\[
 \frac{K_{\rm samp}\eps}{2}\le1-(1-\eps)^{K_{\rm samp}}\le K_{\rm samp}\eps,
 \qquad
 \frac{MK_{\rm samp}\eps}{8}\le\mu(\eps)\le\frac{MK_{\rm samp}\eps}{2}.
\]
Consequently
\begin{equation}\label{eq:q-ratio}
 \frac\nu\mu\le
 \frac{8N_{\rm np}}M+
 8\left(1+\frac{N_{\rm np}}M\right)\frac{J}{K_{\rm samp}}.
\end{equation}
Given an overlay target \(0<\delta_{\rm ov}\le1\), choose
\(M\ge\max(N_{\rm np},32N_{\rm np}/\delta_{\rm ov})\), compile the
finite non-sampling circuit and its \(J\), and then choose
\(K_{\rm samp}\ge128J/\delta_{\rm ov}\). The ratio is at most
\(3\delta_{\rm ov}/8\). At \(\eps_0=1/(2MK_{\rm samp})\), the refresh rate is at
least \(1/16\), and it is strictly increasing from zero. Thus every
nominal interval below \(1/100\) has a physical preimage strictly
below \(\eps_0\). All rate comparisons in the infinite list stay in
the proven range of~\eqref{eq:q-ratio}.

An intrinsic test sees only \(O(W)\) site/round source names in fixed characteristic
translates. Since \(W\mu\le C\), its overlay probability is at most
\(c_{\rm ov}C\delta_{\rm ov}\), with a fixed support count
\(c_{\rm ov}\). The overlap colours and old/new attribution are
unchanged: sampling and refresh add only on-site depth. This replaces
\eqref{eq:overlay} in~\eqref{eq:gamma}. Also
\(p=\Pr(A\lor E)\le(1+\delta_{\rm ov})\mu\le2\mu\).
The \(\mathrm Q\) site has \(M+N_{\rm np}\) wires, all noisy after every layer.
Arbitrary initial probe, helper and scratch states are reset before
use. Every outside data slot retains its own selector gatekeeper.
Thus neither the new hardware nor the refresh can bypass the original
per-round complete cut.
\subsection{A terminating choice of all constants}\label{s:parameters}
The choices specify a fixed medium by terminating algorithms. First
fix the real semantic gate inventory, roles, autonomy and complete-cut
templates. The semantic width, halo and participation bounds determine
\(a_{\rm sep}\) and \(b_{\rm FT}\) in \eqref{d:budgets}.
Choose \(e=8\), \(r_{\rm FT}=8b_{\rm FT}\), the least code parameter
in \eqref{d:capacity}, the hardened finite quantum table, and then
\(P_{\rm cal},T_{\rm cal}\) and its classical interpreter.
This fixes the full witness geometry and entropy used by the positive
threshold. Physical conversion multiplicity has already been paid
in \(b_{\rm FT}\).

Choose rational \(0<C<.01\) small enough for the strict source
inequalities. Search powers \(B=4^k\), with \(B\ge20/C\), until
\(4z_B<1\) and every intrinsic source base is below its required
bound. The quantities in \eqref{eq:c-eta} decay exponentially in
\(B\), while their support and entropy prefactors grow polynomially.
This search terminates. With \(s_B=\lceil3CB/4\rceil\), it also gives
\(s_B/B\le.8C\), below the long-helper mean.

After \(B\), choose support colours, overlay margins and a positive
macro-rate cap. Decrease it to meet the source-certificate recursion
of Lemma~\ref{d:recursion}, the historical/boundary and rare-path
bounds of Section~\ref{sec:controller}, and the wrapped cap of
Lemma~\ref{u:wrapped}. Deep queries require at least sixteen raw names:
their lag entropy \(W^{6f/16}\) is beaten by raw weight
\((2\mu)^{3f/4}\), leaving the positive power \(\mu^{3f/8}\).
Together with \(W\mu\le C\), this makes the smallness searches
terminate. The cap is chosen for this fixed processor catalogue.

Choose \(\mu_*\le\min(C/100,1/100,p_*/2)\) inside all these margins
and then \(n_0\ge k+2\) with
\(21C/[10\ell_{n_0+1}]<\mu_*\). Build the fixed scan machine,
its chart alphabet and binary implementation only now. Full source
geometry and cap bounds do not depend on the number of palette codes.
In \(\mathrm W\), count every actual internal layer \(J\) and choose
independent padding \(P_{\rm pad}\) large enough for
\eqref{eq:overlay}. In \(\mathrm Q\), choose the overlay target,
count all nonprobe wires, choose \(M\), count the actual nonsampling
\(J\), and then choose \(K_{\rm samp}\) as above. Every port, idle,
reset and decomposition layer is included. Long on-site depth changes
these counts without adding a spatial quantum stencil.

For \(n\ge0\), define
\begin{equation}\label{eq:windows}
 J_n=\left(\frac{2C}{\ell_{n+n_0+1}},
            \frac{21C}{10\ell_{n+n_0+1}}\right).
\end{equation}
Use the separately chosen strictly increasing rational refresh polynomial
\(\mu_{\mathrm W}\) or \(\mu_{\mathrm Q}\). Enumerate rational pairs
\(0<a_n<b_n\) in the overlay-valid regime until
\[
 \inf J_n<\mu(a_n)<\mu(b_n)<\sup J_n,
 \qquad I_n=(a_n,b_n).
\]
Exact rational polynomial comparisons suffice, and density guarantees
termination. These are rational inner intervals; their closures lie
strictly inside the full preimages, whose endpoints can be algebraic.
They are nonempty, disjoint and tend to zero. The offset calculation
following Proposition~\ref{c:selector-import} identifies their exact
machine input. Choose the rational storage coefficient \(c\) after
the final speed and source footprint, and the total threshold \(L_0(n)\)
after all finite label margins, as in Proposition~\ref{d:storage}.
